% ============================================================================
%  VSEPR-SIM — §L: Five-Layer Vertical Integration Stack
%  Materials Discovery Engine.  Full Fuel Cycle Simulation.
%  Cradle to Product.
%  Printer-ready LaTeX.  Compile with pdflatex (twice for refs).
% ============================================================================
\documentclass[11pt, a4paper, oneside]{article}

% ── Geometry ────────────────────────────────────────────────────────────────
\usepackage[
  top=1.0in, bottom=1.0in, left=1.15in, right=1.15in,
  headheight=14pt
]{geometry}

% ── Fonts / encoding ───────────────────────────────────────────────────────
\usepackage[T1]{fontenc}
\usepackage{lmodern}
\usepackage{microtype}

% ── Math ───────────────────────────────────────────────────────────────────
\usepackage{amsmath, amssymb, amsthm, mathtools}
\usepackage{bm}

% ── Tables ─────────────────────────────────────────────────────────────────
\usepackage{booktabs}
\usepackage{array}
\usepackage{tabularx}
\newcolumntype{L}[1]{>{\raggedright\arraybackslash}p{#1}}
\newcolumntype{C}[1]{>{\centering\arraybackslash}p{#1}}

% ── Lists ──────────────────────────────────────────────────────────────────
\usepackage{enumitem}
\setlist{nosep, leftmargin=1.5em}

% ── Code listings ─────────────────────────────────────────────────────────
\usepackage{listings}
\usepackage{xcolor}
\definecolor{lstbg}{RGB}{248,248,248}
\definecolor{lstkw}{RGB}{0,0,180}
\definecolor{lstcm}{RGB}{80,120,80}
\lstdefinestyle{cppstyle}{
  language=C++, basicstyle=\ttfamily\footnotesize,
  backgroundcolor=\color{lstbg},
  keywordstyle=\color{lstkw}\bfseries,
  commentstyle=\color{lstcm}\itshape,
  frame=single, framesep=3pt, breaklines=true, keepspaces=true, tabsize=4,
  showstringspaces=false
}
\lstset{style=cppstyle}

% ── Theorem environments ───────────────────────────────────────────────────
\theoremstyle{definition}
\newtheorem{definition}{Definition}[section]
\newtheorem{postulate}[definition]{Postulate}
\newtheorem{remark}[definition]{Remark}
\newtheorem{hardproblem}[definition]{Hard Problem}

% ── TColorbox ─────────────────────────────────────────────────────────────
\usepackage{tcolorbox}
\tcbuselibrary{skins, breakable}

\definecolor{l1col}{RGB}{230,245,255}
\definecolor{l2col}{RGB}{255,245,220}
\definecolor{l3col}{RGB}{240,255,230}
\definecolor{l4col}{RGB}{255,230,240}
\definecolor{l5col}{RGB}{235,225,255}
\definecolor{hardcol}{RGB}{255,220,210}

% ── Headers / footers ─────────────────────────────────────────────────────
\usepackage{fancyhdr}
\pagestyle{fancy}
\fancyhf{}
\fancyhead[L]{\small\textit{VSEPR-SIM Research Platform}}
\fancyhead[R]{\small\textit{§L — Five-Layer Vertical Integration Stack}}
\fancyfoot[C]{\thepage}
\renewcommand{\headrulewidth}{0.4pt}

% ── Section formatting ─────────────────────────────────────────────────────
\usepackage{titlesec}
\titleformat{\section}{\Large\bfseries}{§\thesection}{1em}{}
\titleformat{\subsection}{\large\bfseries}{\thesubsection}{1em}{}
\titleformat{\subsubsection}{\normalsize\bfseries}{\thesubsubsection}{1em}{}

% ── Hyperlinks ─────────────────────────────────────────────────────────────
\usepackage{hyperref}
\hypersetup{colorlinks=true,linkcolor=blue!60!black,urlcolor=blue!70!black}

% ── Macros ─────────────────────────────────────────────────────────────────
\newcommand{\file}[1]{\texttt{#1}}
\newcommand{\type}[1]{\texttt{#1}}
\newcommand{\Ang}{\text{\AA}}
\newcommand{\kcalmol}{\text{kcal/mol}}
\newcommand{\boldI}{\bm{I}}
\newcommand{\lname}[1]{\textbf{L#1}}

% ============================================================================
\begin{document}
% ============================================================================

\begin{titlepage}
  \centering
  \vspace*{1.8cm}
  {\Huge\bfseries VSEPR-SIM Research Platform\par}
  \vspace{0.8cm}
  {\Large\bfseries §L\quad Five-Layer Vertical Integration Stack\par}
  \vspace{0.5cm}
  \rule{0.7\textwidth}{0.6pt}\par
  \vspace{0.8cm}
  {\large\itshape
    Materials Discovery Engine\\[0.3em]
    Full Fuel Cycle Simulation\\[0.3em]
    Cradle to Product\par}
  \vspace{1.2cm}
  \begin{tabular}{cl}
    \toprule
    Layer & Description \\
    \midrule
    L1 & Paper Identity — Formula, class, name.  Static. \\
    L2 & Atomistic — Z, A, Q, $\varepsilon$.  Physics and energy state. \\
    L3 & Molecular 3D — Geometry, bonding, coordination. \\
    L4 & Atomistic / CG — Beads, phase behaviour, bulk interactions. \\
    L5 & Macro — CAD-scale geometry, volumes, surfaces. \\
    \bottomrule
  \end{tabular}
  \vspace{1.2cm}

  {\normalsize
    The particle identity from L2 propagates all the way up.\\
    The iron in a pipe at L5 is still Fe with its $Z$, $A$, $Q$ underneath.\\
    The macro geometry is the accumulated consequence of everything below it.\par
  }
  \vfill
  {\small
    Implementation files:\quad
    \file{include/layer\_stack.hpp} \quad
    \file{coarse\_grain/models/layer\_boundary.hpp}\par
    \vspace{0.3em}
    Status: Research-oriented active development
  }
\end{titlepage}

\tableofcontents
\clearpage

% ============================================================================
\section{Architecture Overview}
\label{sec:overview}
% ============================================================================

\subsection{Purpose}

The five-layer stack is the routing document for the entire VSEPR-SIM
platform.  It answers the question: \emph{what is the particle at every
scale, and how does it get there?}

The stack is not a layer model in the OSI sense — layers do not replace
each other.  They accumulate.  Each higher layer is an interpretation of
the same underlying physical object at a different resolution.

\subsection{The Vertical Integration Invariant}

\begin{postulate}[Vertical Integration Invariant]
For any particle $i$ and any layer $k \in \{2, 3, 4, 5\}$:
\[
  L_k(i).Z = L_2(i).Z, \qquad L_k(i).A = L_2(i).A
\]
Nuclear identity is immutable across all layers under all non-nuclear
transformations.  The iron in a pipe at L5 is still Fe at L2.
\end{postulate}

Charge $Q_i$ and energy state $\varepsilon_i$ are \emph{environment-dependent}
and evolve as the chemical context changes:
\[
  \mathrm{Fe}^0\;(\text{ore, L2}) \;\to\;
  \mathrm{Fe}^{3+}\;(\text{acid leach, L2}) \;\to\;
  \mathrm{Fe}_2\mathrm{O}_3\;\text{bead}\;(\text{L4})
\]
Same $Z = 26$.  Different $Q$, different $\varepsilon$, different $\Sigma_i$
(structural role changes from metallic to ionic at oxidation).

\subsection{The Compute Problem}

The discovery-end use case requires pushing millions of candidates through
the stack.  Most die at L1.  The rest are screened progressively:

\begin{center}
\begin{tabular}{lrll}
\toprule
Layer & Typical survivors & Cost per candidate & Gate criterion \\
\midrule
L1 screen & $10^6$ & $\mu$s & Formula valid, class not excluded \\
L2 screen & $10^4$ & ms & $\varepsilon$ in range, $|Q| < Q_{\max}$, $\Lambda \geq \Lambda_{\min}$ \\
L3 screen & $10^3$ & 0.1--1 s & Geometry converged, $F_{\rm rms} < F_{\rm tol}$ \\
L4 screen & $10^2$ & 1--10 s & Bead mapping residual $< \delta$, phase coherent \\
L5 report & $10^1$ & 10--100 s & CAD geometry, corrosion model, integrity check \\
\bottomrule
\end{tabular}
\end{center}

This is the Hourglass model applied at the platform level (§R).  Each
layer is a gate.  The funnel is by design.

% ============================================================================
\section{L1 — Paper Identity}
\label{sec:l1}
% ============================================================================

\begin{tcolorbox}[colback=l1col, colframe=blue!40!black,
  title={L1 — Paper Identity},fonttitle=\bfseries\small]
  Formula, class, name.  Static.  Never mutates.\\
  A candidate that fails L1 is never instantiated at L2 or above.
\end{tcolorbox}

\subsection{What L1 carries}

\begin{itemize}
  \item \textbf{Formula} — Hill-order chemical formula string
        (\texttt{"Fe2O3"}, \texttt{"UO2"}, \texttt{"Ti6Al4V"})
  \item \textbf{IUPAC name} — systematic name (may be empty for elements)
  \item \textbf{Common name} — engineering or mineral name
        (\texttt{"hematite"}, \texttt{"urania"}, \texttt{"sphalerite"})
  \item \textbf{CAS number} — registry identifier
  \item \textbf{Class tag} — coarse classification:
        \texttt{"element"}, \texttt{"oxide"}, \texttt{"halide"},
        \texttt{"organometallic"}, \texttt{"alloy"}, \texttt{"mineral"}
  \item \textbf{$Z$-set and stoichiometry} — atomic numbers and coefficients
\end{itemize}

\subsection{L1 gate}

A species passes L1 if and only if:
\begin{enumerate}
  \item Formula is non-empty and parseable
  \item All $Z_i \in [1, 118]$
  \item \texttt{class\_tag} is not in the exclusion list
\end{enumerate}

L1 costs microseconds.  It is the first and cheapest filter.  Run it on
the full candidate set before doing anything else.

\subsection{Process simulation context}

In a full fuel cycle workflow, L1 represents the catalog:
\begin{itemize}
  \item Every ore mineral has an L1 record (\texttt{"Fe2O3"},
        \texttt{"FeS2"}, \texttt{"FeCO3"})
  \item Every reagent has an L1 record (\texttt{"H2SO4"},
        \texttt{"NaOH"}, \texttt{"HNO3"})
  \item Every product target has an L1 record (\texttt{"Fe"},
        \texttt{"UO2"}, \texttt{"ThO2"})
\end{itemize}

The L1 catalog is populated once and queried throughout the simulation.
No physics is required at this stage.

% ============================================================================
\section{L2 — Atomistic Layer}
\label{sec:l2}
% ============================================================================

\begin{tcolorbox}[colback=l2col, colframe=orange!60!black,
  title={L2 — Atomistic},fonttitle=\bfseries\small]
  $Z$, $A$, $Q$, $\varepsilon$.  Physics and energy state.\\
  Environment-dependent.  The §0 Identity Vector lives here.
\end{tcolorbox}

\subsection{What L2 carries}

\begin{center}
\begin{tabular}{llll}
\toprule
Field & Symbol & Units & Notes \\
\midrule
Atomic number & $Z_i$ & — & Immutable. Anchor of propagation invariant. \\
Mass number   & $A_i$ & — & Immutable. 0 = natural-abundance average. \\
Charge        & $Q_i$ & $e$ & Environment-dependent. $\mathrm{Fe}^{3+}$ vs $\mathrm{Fe}^0$. \\
LJ well depth & $\varepsilon_i$ & kcal/mol & Energy state. Scales with polarisability. \\
LJ radius     & $\sigma_i$ & \Ang & Size parameter. \\
Temperature   & $T_i$ & K & Local thermal state. \\
Pressure      & $P_i$ & bar & Local mechanical state. \\
Phase         & $\phi_i$ & — & \texttt{SOLID, LIQUID, GAS, AQUEOUS, \ldots} \\
\midrule
$\Sigma_i$ & StructuralRole & — & Bonding topology prior (§0). \\
$\Lambda_i$ & StabilityClass & — & Persistence class (§0). \\
\bottomrule
\end{tabular}
\end{center}

\subsection{The environment-dependence of Q and $\varepsilon$}

This is not an approximation — it is physics.
$\mathrm{Fe}^{3+}$ in sulphuric acid leach solution behaves fundamentally
differently from $\mathrm{Fe}^0$ in a solid ore body:

\begin{equation}
  \mathrm{Fe}^0 \;\xrightarrow{\text{oxidation}}\; \mathrm{Fe}^{3+}
  \qquad
  Q_i: 0 \;\to\; +3
  \quad
  \Sigma_i: \text{Metallic} \;\to\; \text{Ionic-Dominant}
  \label{eq:fe-oxidation}
\end{equation}

The $\varepsilon$ track this transition:
\begin{itemize}
  \item $\mathrm{Fe}^0$: $\varepsilon \approx 0.13$\,kcal/mol (metallic, dispersion-dominant)
  \item $\mathrm{Fe}^{3+}$ (aq): $\varepsilon \approx 0.55$\,kcal/mol (ionic, Coulomb-dominant)
\end{itemize}

Your $\varepsilon$ field tracks that energy transition at every simulation step.

\subsection{L2 gate}

Candidates fail L2 if:
\begin{itemize}
  \item $\varepsilon < 0$ or $\varepsilon > \varepsilon_{\max}$ (unphysical energy state)
  \item $|Q_i| > Q_{\max}$ (exceeds ionic saturation)
  \item $\Lambda_i < \Lambda_{\min}$ (unstable in this environment)
\end{itemize}

% ============================================================================
\section{L3 — Molecular 3D Layer}
\label{sec:l3}
% ============================================================================

\begin{tcolorbox}[colback=l3col, colframe=green!50!black,
  title={L3 — Molecular 3D},fonttitle=\bfseries\small]
  Geometry, bonding, coordination, local environment.\\
  Built by FIRE minimisation on the L2 state.
\end{tcolorbox}

\subsection{What L3 carries}

\begin{itemize}
  \item Optimised 3D positions (Å) for all atoms in the species
  \item Bond table: pairs $(i,j)$, order, length, bond energy proxy
  \item Angle table: triples $(i,j,k)$ with deviation from ideal
  \item Coordination record per atom: coordination number, geometry class,
        bond-length spread (distortion indicator)
  \item Quality metrics: $F_{\rm rms}$, total energy, convergence flag
  \item Local environment summary: mean coordination, anisotropy index
\end{itemize}

\subsection{L3 in the process simulation}

The acid leach stage is an L3 event.  When
$\mathrm{Fe}_2\mathrm{O}_3$ dissolves in $\mathrm{H}_2\mathrm{SO}_4$:
\begin{enumerate}
  \item L1: \texttt{"Fe2O3"} + \texttt{"H2SO4"} are both in the catalog
  \item L2: $Q(\mathrm{Fe})$ increments from $0 \to +3$ over the dissolution timestep
  \item L3: Fe–O bonds break; Fe–O(H$_2$O) coordination shell forms
  \item L4: The solvated Fe$^{3+}$ complex becomes a charged CG bead
\end{enumerate}

L3 is where the bonding topology changes are made explicit.

% ============================================================================
\section{L4 — Atomistic Bead Layer}
\label{sec:l4}
% ============================================================================

\begin{tcolorbox}[colback=l4col, colframe=red!50!black,
  title={L4 — Atomistic Bead (CG)},fonttitle=\bfseries\small]
  Coarse-grained beads, phase behaviour, bulk interactions.\\
  Anisotropic surface mapping lives here.
\end{tcolorbox}

\subsection{What L4 carries}

\begin{itemize}
  \item The full \type{Bead} struct from the CG layer, including:
    \begin{itemize}
      \item $\Sigma_i$ (structural role) and $\Lambda_i$ (stability class)
      \item Three-channel role-weighted interaction kernel weights
      \item Anisotropic surface descriptor (Fuzzy Ball Model)
      \item Slow environmental state $\eta$ from the 6+9 Stepper
    \end{itemize}
  \item Phase region: dominant phase + volume fraction + order parameter $P_2$
  \item Local density $\rho_B$ and coordination number $C_B$
  \item Back-reference to the dominant L2 atom (preserving $Z$, $A$, $Q$)
\end{itemize}

\subsection{Anisotropic surface mapping}

The anisotropic surface descriptor is the Fuzzy Ball Model (§Fuzzy).
At L4, the bead's surface carries a three-channel spherical harmonic
field:
\[
  S^{(k)}(\theta, \phi) = \sum_{\ell=0}^{\ell_{\max}}
                           \sum_{m=-\ell}^{+\ell}
                           c_{\ell m}^{(k)} Y_\ell^m(\theta, \phi),
  \quad k \in \{\text{steric}, \text{electrostatic}, \text{dispersion}\}
\]
This is the Container Geometry Approximation $\Omega_i$ from §0,
made concrete.

% ============================================================================
\section{L5 — Macro / CAD Layer}
\label{sec:l5}
% ============================================================================

\begin{tcolorbox}[colback=l5col, colframe=violet!60!black,
  title={L5 — Macro / CAD},fonttitle=\bfseries\small]
  Actual geometry, volumes, surfaces, material bodies.\\
  A pipe. A reactor wall. A valve seat.
\end{tcolorbox}

\subsection{What L5 carries}

\begin{center}
\begin{tabular}{lll}
\toprule
Field & Units & Description \\
\midrule
\texttt{body\_type} & — & Pipe, Vessel, Valve, HeatExchanger, \ldots \\
\texttt{material\_formula} & — & L1 formula of bulk material \\
\texttt{length\_m} & m & Body length \\
\texttt{outer\_diameter\_m} & m & Outer diameter \\
\texttt{wall\_thickness\_m} & m & Wall thickness \\
\texttt{volume\_m3} & m$^3$ & Total enclosed volume \\
\texttt{surface\_area\_m2} & m$^2$ & Wetted or exposed surface area \\
\midrule
\texttt{corrosion\_depth\_m} & m & Wall loss from chemical attack \\
\texttt{oxide\_layer\_m} & m & Passivation / scale layer thickness \\
\texttt{roughness\_Ra\_m} & m & Surface roughness \\
\texttt{dominant\_species} & — & L1 formula of the surface species \\
\midrule
\texttt{temperature\_K} & K & Operating temperature \\
\texttt{pressure\_bar} & bar & Operating pressure \\
\texttt{fluid\_formula} & — & L1 formula of the process fluid \\
\texttt{structural\_integrity} & 0--1 & 1 = pristine, 0 = failed \\
\bottomrule
\end{tabular}
\end{center}

\subsection{Particle identity at L5}

The iron in a pipe wall at L5 is still Fe ($Z = 26$) underneath.
The macro geometry — wall thickness, corrosion depth, surface roughness —
is the accumulated consequence of:
\begin{itemize}
  \item L2: Fe$^0$ $\varepsilon$ determining surface adhesion energy
  \item L3: Fe–O bond formation rate at the surface
  \item L4: oxide scale growth modelled by bead density and phase fields
  \item L5: wall loss recorded as \texttt{corrosion\_depth\_m}
\end{itemize}

The structural integrity field is a direct function of these upward-propagated
identities:
\[
  \text{integrity} = 1 - \frac{\delta_{\rm corr}}{t_{\rm wall}}
  \quad \in [0, 1]
\]
where $\delta_{\rm corr}$ is the corrosion depth and $t_{\rm wall}$ is the
original wall thickness.

% ============================================================================
\section{The L2 ↔ L4 Hard Boundary}
\label{sec:boundary}
% ============================================================================

\begin{tcolorbox}[colback=hardcol, colframe=red!70!black,
  title={Hard Problem 1: The L2 ↔ L4 Phase Transition Boundary},
  fonttitle=\bfseries\small, breakable]
This is where most simulations cheat or break.\\[4pt]
Atomic behaviour (discrete particle physics, quantum-derived $\varepsilon$,
explicit bonding) transitions into coarse-grained bulk behaviour (phase
fields, density order parameters, collective interactions).
It is not a sharp line.  It is a controlled approximation zone.
\end{tcolorbox}

\subsection{What cheating looks like}

\begin{enumerate}
  \item \textbf{Silent $\varepsilon$ switch.}
        Swapping from LJ $\varepsilon$ (L2) to CG bead $\varepsilon$ (L4)
        without recording the mapping error.  The energy changes by an
        unmeasured amount.

  \item \textbf{Charge loss.}
        Lumping Fe$^{3+}$ and O$^{2-}$ into a neutral bead loses the
        Coulomb interaction entirely.  The electrostatic energy disappears
        from the simulation without a trace.

  \item \textbf{Single-phase fiction.}
        Treating a bead group as single-phase when atoms in it straddle
        a liquid-solid boundary.  Phase boundaries do not coincide with
        CG mapping boundaries.

  \item \textbf{Frozen structural role.}
        Using the same $\Sigma_i$ = Metallic for Fe across the
        liquid-metal/ionic transition.  The bonding topology changes;
        the role must follow.
\end{enumerate}

\subsection{The \file{layer\_boundary.hpp} solution}

Five boundary checks run at every L2→L4 transition:

\begin{center}
\begin{tabular}{llp{7cm}}
\toprule
Unit & Name & What it measures \\
\midrule
B1 & Energy consistency &
  $|E_{\rm L4} - E_{\rm L2}|$ per bead group.
  Flags when CG energy deviates from atomistic. \\
B2 & Charge conservation &
  $|q_{\rm L4} - \sum_j q_j^{\rm L2}|$ per bead.
  Non-conservation above threshold is always flagged. \\
B3 & Phase coherence &
  Detects bead groups spanning multiple phases.
  Optional split or flag. \\
B4 & Structural role re-evaluation &
  Re-classifies $\Sigma_i$ from current $Q_i$ and $Z_i$.
  Updates the bead in-place if changed. \\
B5 & Mapping residual &
  Weighted composite $(0.5 \times E_{\rm term}
    + 0.3 \times Q_{\rm term}
    + 0.2 \times \phi_{\rm term})$.
  High residual $>$ 0.15 is flagged. \\
\bottomrule
\end{tabular}
\end{center}

The boundary report is a first-class output.  Every approximation is
recorded with its magnitude.  Nothing is silent.

\subsection{The B4 role re-evaluation equation}

\begin{equation}
  \Sigma_i^{(\rm boundary)} =
  \begin{cases}
    \text{IonicDominant} & Z_i \in \text{transition metals}
                           \;\wedge\; |Q_i| > 0.8 \\
    \text{Metallic}      & Z_i \in \text{transition metals}
                           \;\wedge\; |Q_i| \leq 0.8 \\
    \text{IonicDominant} & |Q_i| > 0.5 \\
    \text{DirectionalCovalent} & \text{otherwise}
  \end{cases}
  \label{eq:b4-role}
\end{equation}

This is equation~(\ref{eq:fe-oxidation}) made computational.

% ============================================================================
\section{Hard Problem 2: Compute Cost}
\label{sec:compute}
% ============================================================================

\begin{tcolorbox}[colback=hardcol, colframe=red!70!black,
  title={Hard Problem 2: Compute Cost},
  fonttitle=\bfseries\small]
Pushing millions of candidates through all five layers is not feasible
without aggressive early rejection and lazy evaluation.
\end{tcolorbox}

\subsection{The answer: lazy layer evaluation}

No particle is evaluated at layer $k$ until it has passed the gate at
layer $k-1$.  The \type{LayerParticle} struct uses \texttt{std::optional}
for L3, L4, and L5 precisely for this reason:

\begin{lstlisting}
struct LayerParticle {
    L1_PaperIdentity                     l1;  // always populated
    L2_AtomisticState                    l2;  // populated after L1 gate
    std::optional<L3_MolecularGeometry>  l3;  // populated after L2 gate
    std::optional<L4_AtomisticBeadState> l4;  // populated after L3 gate
    std::optional<L5_MacroGeometry>      l5;  // populated after L4 gate
};
\end{lstlisting}

L3, L4, and L5 are never heap-allocated for particles that fail their
predecessor gate.  The rejected candidate costs only its screen cost,
not the full simulation cost.

\subsection{The answer: Hourglass at platform scale}

The Hourglass Model (§R) was designed for individual configuration
sampling.  At platform scale the same logic applies to the candidate
population:

\begin{enumerate}
  \item \textbf{Mouth:} Generate $10^6$ L1 candidates from combinatorial
        formula space
  \item \textbf{Neck:} Sequential L1/L2/L3 gates reduce to $\sim 10^2$
  \item \textbf{Base:} Surviving candidates get full L4/L5 treatment and
        are ranked by composite score
\end{enumerate}

The gate criteria are the same physics as the layer boundaries.
They are not ad hoc filters — they are the physics screening.

% ============================================================================
\section{Full Fuel Cycle Simulation}
\label{sec:fuel-cycle}
% ============================================================================

\subsection{Discovery end}

Generate candidate metal species, classify, screen.  Millions of
candidates, most rejected at L1, survivors get full L2–L5 treatment.
This document is the routing specification for that pipeline.

\begin{center}
\begin{tabular}{lll}
\toprule
Stage & Input & Output \\
\midrule
Candidate generation & L1 formula space & $10^6$ L1 records \\
L1 screen & Class + formula validity & $\sim 10^4$ survivors \\
L2 evaluation & $\varepsilon$, $Q$, $\Lambda$ screening & $\sim 10^3$ survivors \\
L3 geometry & FIRE minimisation & $\sim 10^2$ converged structures \\
L4 mapping & CG bead + phase + anisotropy & $\sim 10^1$ viable materials \\
L5 report & CAD geometry, corrosion, integrity & Product candidates \\
\bottomrule
\end{tabular}
\end{center}

\subsection{Process simulation}

Ore input $\to$ crushing $\to$ dissolution $\to$ separation $\to$
refining.  Every stage is a chemical environment.  Species identities
evolve as they move through:

\begin{center}
\begin{tabular}{lll}
\toprule
Stage & Active layer(s) & Key identity change \\
\midrule
Ore body & L1, L4 & Fe as \texttt{Fe2O3}, \texttt{FeS2} \\
Crushing / grinding & L5 & Surface area increases \\
Acid leach & L2, L3 & Fe$^{3+}$ forms; $Q$, $\Sigma$ update \\
Solvent extraction & L3, L4 & Complex speciation; bead re-mapping \\
Stripping & L2, L3 & Metal ion concentration \\
Electrowinning & L2, L5 & Metal deposits; cathode grows \\
Refining & L3, L4, L5 & Alloy composition; grain structure \\
\bottomrule
\end{tabular}
\end{center}

\subsection{Pipe and reactor simulation}

L3 to L5.  Material integrity, corrosion, fluid-wall interactions,
thermal loads.  The CAD layer is the physical infrastructure.
The particle layer is what is eating it or flowing through it.

The corrosion loop:
\begin{enumerate}
  \item L2: $\varepsilon(\mathrm{Fe})$ and $Q(\mathrm{Fe})$ at the pipe surface
  \item L3: Fe–O bond formation rate (surface oxidation kinetics)
  \item L4: Oxide scale growth rate (bead phase field, $\eta$ slow state)
  \item L5: \texttt{corrosion\_depth\_m} increments; \texttt{structural\_integrity} decrements
\end{enumerate}

% ============================================================================
\section*{Conclusion}
% ============================================================================

The five-layer stack is the architectural spine of the VSEPR-SIM
materials discovery platform.  Its properties:

\begin{enumerate}
  \item \textbf{Vertically integrated.}
        The same particle identity ($Z$, $A$, $Q$) propagates from
        ore body to finished material without discontinuity.

  \item \textbf{Lazy.}
        Candidates are never evaluated at cost above their gate status.
        The funnel is the architecture.

  \item \textbf{Explicit at every boundary.}
        The L2↔L4 hard boundary is not approximated silently.
        \file{layer\_boundary.hpp} records every energy gap, every
        charge error, every role update, every phase mismatch.

  \item \textbf{Anti-black-box.}
        Every layer transition is a named, typed, inspectable C++ struct.
        Every approximation has a magnitude.  Nothing is hidden.

  \item \textbf{Deterministic.}
        Same candidate + same environment + same parameters
        $\Rightarrow$ same layer traversal $\Rightarrow$ same output.
\end{enumerate}

\end{document}
